\section{서론}

대규모 언어 모델(Large Language Model, LLM)의 발전과 함께, 단일 LLM
호출을 넘어 여러 개의 LLM 에이전트가 협업하는 멀티 에이전트(Multi-Agent)
시스템의 활용이 빠르게 증가하고 있다~\cite{wang2024survey, wu2023autogen}.
문서 생성, 코드 작성, 데이터 분석과 같은 복잡한 작업은 단일 에이전트로
처리하기보다, Planner, Writer, Reviewer, Reviser 등 역할이 분담된 여러
에이전트가 협력할 때 더 높은 품질을 얻을 수 있다.

그러나 멀티 에이전트 시스템에서 \emph{어떤 에이전트를 어떤 순서로, 언제
호출할 것인지}를 결정하는 오케스트레이션(Orchestration) 방식은 시스템의
성능을 좌우하는 핵심 요소임에도 충분히 논의되지 않았다. 오케스트레이션은
크게 두 가지로 구분된다. 첫째는 실행 흐름을 코드로 고정하는 \textbf{코드
기반(Code Orchestration)}이고, 둘째는 다음 동작을 LLM이 스스로 결정하는
\textbf{LLM 기반(LLM Orchestration)}이다~\cite{langgraph2024, googleadk2025}.

기존 연구는 주로 개별 에이전트의 성능, 프롬프트 설계, 도구(Tool) 사용
방법에 초점을 맞추어 왔다~\cite{hong2024metagpt}. 반면, 동일한 멀티
에이전트 환경에서 오케스트레이션 \emph{구조 자체}가 실행 시간, 품질,
비용, 일관성에 미치는 영향을 정량적으로 비교한 연구는 부족하다. 구조
선택은 실무에서 빈번하게 마주치는 결정임에도, 이를 뒷받침할 실험적 근거와
설계 지침이 마련되어 있지 않다.

따라서 본 논문에서는 동일한 멀티 에이전트 시스템 위에서 Code, LLM,
Hybrid 세 가지 오케스트레이션 구조를 구현하고 그 성능을 비교한다. 나아가
실험 결과를 바탕으로 작업 특성에 따라 적합한 오케스트레이션을 선택하기
위한 설계 가이드라인(Design Guideline)을 제안한다. 본 논문의 기여는 다음과
같다.

\begin{itemize}
  \item Code, LLM, Hybrid 오케스트레이션 구조를 동일한 멀티 에이전트
        환경에서 구현하고 비교한다.
  \item 실행 시간, 품질, 토큰 사용량, 일관성 등 다양한 성능 지표를
        정량적으로 분석한다.
  \item 실험 결과를 바탕으로 작업 특성에 따른 오케스트레이션 선택 기준
        (Design Guideline)을 제안한다.
\end{itemize}
